% Figure: multiplication tables mod 5 vs mod 6 --- units highlighted.
% Source for Figure~\ref{fig:units}. TikZ only; regenerable by compiling main.tex.
\begin{tikzpicture}[
  cell/.style={draw=rulegray,minimum width=6mm,minimum height=5mm,font=\scriptsize,inner sep=0.6pt},
  un/.style={cell,fill=accentsoft},
  hd/.style={cell,font=\scriptsize\bfseries},
  lab/.style={font=\footnotesize\sffamily,text=slate},
]
% mod 5 table rows 1..4
\node[lab,anchor=east] at (-0.7,0) {$\bmod\ 5$:};
\foreach \c/\h in {0/{},1/1,2/2,3/3,4/4}
  \node[hd] at (\c*0.62,0) {\h};
\foreach \r in {1,2,3,4}
{
  \node[hd] at (0,-\r*0.55) {\r};
  \foreach \c in {1,2,3,4}
  {
    \pgfmathtruncatemacro{\v}{mod(\r*\c,5)}
    \node[un] at (\c*0.62,-\r*0.55) {\v};
  }
}
\node[lab,anchor=west] at (3.3,-1.1) {every nonzero row is a permutation:};
\node[lab,anchor=west] at (3.3,-1.5) {each has a 1-entry $\Rightarrow$ inverse};
% mod 6 table rows 1..5
\node[lab,anchor=east] at (-0.7,-3.0) {$\bmod\ 6$:};
\foreach \c/\h in {0/{},1/1,2/2,3/3,4/4,5/5}
  \node[hd] at (\c*0.62,-3.0) {\h};
\foreach \r in {1,...,5}
{
  \node[hd] at (0,-3.0-\r*0.55) {\r};
  \foreach \c in {1,...,5}
  {
    \pgfmathtruncatemacro{\v}{mod(\r*\c,6)}
    \pgfmathtruncatemacro{\g}{gcd(\r,6)}
    \ifnum\g=1
      \node[un] at (\c*0.62,-3.0-\r*0.55) {\v};
    \else
      \node[cell] at (\c*0.62,-3.0-\r*0.55) {\v};
    \fi
  }
}
\node[lab,anchor=west] at (4.4,-3.8) {shaded rows ($\gcd=1$): still};
\node[lab,anchor=west] at (4.4,-4.2) {permutations $\Rightarrow$ inverses.};
\node[lab,anchor=west] at (4.4,-4.8) {unshaded rows collapse:};
\node[lab,anchor=west] at (4.4,-5.2) {no 1-entry, and $0$ appears ---};
\node[lab,anchor=west] at (4.4,-5.6) {$2\cdot3\equiv0$ kills cancellation.};
\end{tikzpicture}
